\section{Related Work}
\textbf{Causal Discovery from Time Series data.} A prominent approach to time series causal discovery is based on Granger causality \citep{granger1969investigating}, and its extension using neural networks \citep{khanna2019economy, tank2021neural, lowe2022amortized, cheng2023cutsneuralcausaldiscovery, cheng2024cuts+}.  However, Granger causality only captures predictive relationships and ignores instantaneous effects, latent confounders, and history-dependent noise \citep{peters2017elements}. The Structural Causal Model (SCM) framework, as implemented in  \citet{hyvarinen2010estimation, pamfil2020dynotears, gong2022rhino, wang2024neuralstructurelearningstochastic}, can theoretically overcome these limitations by explicitly modeling causal relationships between variables. Another line of work \citep{runge2019detecting, runge2020discovering}  extends the conditional independence testing-based PC \citep{spirtes2000causation} to time series. However, these methods face scaling and accuracy challenges when applied to spatiotemporal data due to high dimensionality and spatial correlation effects that can mask true causal relationships between distant locations  \citep{xavi2022teleconnection}. 


\textbf{Causal Representation Learning.} The primary objective of causal representation learning is to extract high-level causal variables and their relationships from temporal data. \citet{lippe2022citris, lippecausal} explore this topic specifically in the context of interventional time series data. Meanwhile, \citet{yao2021learning, yao2022temporally, chencaring, moriokacausal, li2024identification, lachapelle2024nonparametric} establish identifiability conditions for reconstructing latent time series from observational data. However, their theorems  rely on various assumptions about the underlying latent dynamics such as no instantaneous effects, sufficient variability or sparsity. These conditions can be challenging to validate in practice. They also do not consider the spatial structure critical to spatiotemporal causal discovery. \change{Another line of research is Dynamic Causal Modeling \citep{FRISTON20031273DCM, STEPHAN20103099DCM, FRISTON2013172DCM, Friston2022DCM}, which infers causal relationships in dynamical systems and has been applied to neuroscience. However, DCM assumes that the parameters of the forward model (i.e., the relationship between the latent and observed variables) are known \emph{a priori}.} In contrast, SPACY guarantees identifiability of latent time series from spatiotemporal data with minimal assumptions about the latent-generating process. 

 \textbf{Spatiotemporal Causal Discovery.} Numerous studies have extended Granger causality to spatiotemporal settings, particularly in climate science \citep{GrangerCausalityofCoupledClimateProcessesOceanFeedbackontheNorthAtlanticOscillation, lozano2009grangerspatial, Grangercausalitycarbondioxide}. Another approach to spatiotemporal causal discovery is to perform dimensionality reduction to obtain a smaller number of latent time series and then infer a causal graph among the latent variables \citep{xavi2022teleconnection, Falasca_2024}. A key limitation of these methods is that dimensionality reduction occurs independently of the causal structure in the data. Consequently, the latent variables may not correspond to causally relevant entities. Some studies like \citet{sheth2022stcd} address specific spatial dependencies relevant to the considered domain. \citet{zhao2023generativecausal} extend \citet{yao2022temporally} by incorporating spatial structures using graph convolutional networks. \citet{brouillard2024causal, boussard2023towards} adopt the single-parent assumption, where each observed variable is influenced by only one latent variable. On the other hand, SPACY allows multiple latent parents per observed variable.

% 


% To be added \textcolor{red}{Another line of research is Dynamic causal modeling\citep{FRISTON20031273DCM, STEPHAN20103099DCM, FRISTON2013172DCM, Friston2022DCM}, which infer causal relationships in a dynamical system and assumes that the parameters of the forward model (i.e. the relationship between latent and observed model are know apriori.}



% \section{Related Work}
% In this section, we provide an overview of the literature on causal discovery from time series, causal representation learning and spatiotemporal causal discovery.

% \paragraph{Causal Discovery from time series data.}
% A prominent line of research in time series causal discovery is based on Granger causality, as introduced by \citet{granger1969investigating}. For example, \citet{tank2021neural} use component-wise MLPs and LSTMs with sparsity constraints to infer non-linear Granger causality. In contrast, \citet{khanna2019economy} apply Statistical Recurrent Units (SRUs) to detect causal relationships across multiple scales. \citet{lowe2022amortized} propose Amortized Causal Discovery using a variational autoencoder and Graph Neural Networks. \citet{cheng2023cutsneuralcausaldiscovery, cheng2024cuts+} infers Granger causal links from irregularly sampled or incomplete data by simultaneously imputing missing values.  However, Granger causality only captures predictive relationships and ignores instantaneous effects, latent confounders, and history-dependent noise \citep{peters2017elements}.

% The Structural Causal Model (SCM) framework can theoretically overcome these limitations by explicitly modeling the causal relationships between variables. \citet{hyvarinen2010estimation} extend LiNGAM \citep{shimizu2006linear} to develop VARLiNGAM, incorporating vector autoregressive models for time series data. DYNOTEARS \citep{pamfil2020dynotears} adapts NOTEARS \citep{zheng2018dags} to dynamic Bayesian networks. Both methods, however, are restricted to linear relationships. PCMCI and PCMCI$^+$ \citep{runge2019detecting, runge2020discovering} extend the PC \citep{spirtes2000causation} algorithm to handle instantaneous effects. Rhino \citep{gong2022rhino} uses neural networks to model the functional relationships and estimates the temporal adjacency matrix from observational data while accounting for exogenous history-dependent noise distributions. \citet{wang2024neuralstructurelearningstochastic} use stochastic differential equations for causal structure learning from continuous-time temporal processes with potentially irregular sampling.

% However, applying these methods directly to spatiotemporal data presents significant challenges. The high dimensionality of large gridded datasets makes it difficult for many techniques—especially those that rely on conditional independence tests—to scale effectively \citep{glymour2019review}. Furthermore, spatially proximate points often exhibit highly correlated and redundant time series. Conditioning on these nearby correlated points can obscure true causal relationships between distant locations, reducing statistical power and leading to inaccurate results \citep{xavi2022teleconnection}.


% \paragraph{Spatiotemporal Causal Discovery/Causal Representation Learning}

% Numerous studies have extended Granger causality to spatiotemporal settings, particularly in climate science \citep{GrangerCausalityofCoupledClimateProcessesOceanFeedbackontheNorthAtlanticOscillation, Grangercausalitycarbondioxide, ali2024causalityearthscience}. \citet{lozano2009grangerspatial} proposed a method combining Granger causality with a Group Elastic Net to capture spatial and temporal dependencies, enabling the identification of causal relationships among climate variables. However, the model assumes that the causal relationships are linear, and only infers a summary graph.

% One approach to spatiotemporal causal discovery is to perform dimensionality reduction to obtain a smaller number of latent time series and then infer a causal graph among the latent variables. For example, \citet{xavi2022teleconnection} use Varimax for dimensionality reduction and PCMCI$^+$ \citep{runge2018pcmci, runge2020pcmci+} for causal discovery. \citet{Falasca_2024} infer regional modes based on correlation and spatial proximity, applying linear-response theory to uncover causal links. A key limitation of these methods is that dimensionality reduction occurs independently of the causal structure in the data. Consequently, the latent variables may not correspond to causally relevant entities. Additionally, conditional independence-based methods are computationally intensive as they may require an exponential number of conditional independence tests.

% Other approaches use neural networks to model nonlinear interactions. The Spatial-Temporal Causal Discovery Framework (STCD) \citep{sheth2022stcd} utilizes attention-based convolutional neural networks to identify causal relationships from gridded time series data. However, it encodes an explicit form of spatial dependence specific to the problem of hydrological systems (i.e. reduce attention scores based on geographic height) rather than inferring it from data. 

% Causal representation learning from time series involves inferring abstract, high-level causal variables and their relationships from temporal data. \citet{lippe2022citris, lippecausal} focus on causal representation learning from interventional time series data. \citet{yao2021learning, yao2022temporally, chencaring} introduce frameworks to recover latent causal variables and identify their relations from observational sequential data. However, these methods do not model instantaneous edges in the causal graph. \citet{moriokacausal} prove the identifiability of causal relations, even in the presence of instantaneous edges, by assuming that the observational variables can be appropriately grouped. However, this grouping is rarely known in practice apriori. 

% Moreover, none of these methods consider the spatial structure present in the data. \citet{zhao2023generativecausal} extends upon \citet{yao2022temporally} by encoding the spatial structure using graph convolutional network. The work most closely related to ours is Causal Discovery with Single Parent Decoding (CDSD) \citep{brouillardcausal, boussard2023towards}, which learns a mapping from the observational time series to latent variables. However, CDSD operates under the assumption that each observed variable is influenced by only one latent variable, and the causal graph has no instantaneous edges. In contrast, SPACY allows both instantaneous edges and overlapping modes, i.e., an observational variable can be influenced by more than one latent variable.

